% generated from Bio 6/rc_double_cover.tex
\chapter{Double strands without a new lemma: the reverse-complement double cover of a de Bruijn graph}\label{chap:Bio6}
\chaptermark{Bio6}
\begin{chapterabstract}
\noindent
The single-strand theory of the companion notes sees a repeat only between co-oriented
copies, so the seven ribosomal operons of \emph{E.\ coli} --- five on one strand, two on
the other --- appear as a five-copy family and a two-copy family and never as a
seven-copy one. Repairing this was expected to need a sandpile theory for bidirected
graphs, which does not exist. It does not need one. Let $D_k(S)$ be the multidigraph
whose arcs are the $(k{+}1)$-mer occurrences in $S$ \emph{and} in its reverse complement
$\bar S$. Then $\mult_D(v)=\mult_S(v)+\mult_S(\bar v)$ and
$d^{+}(v)=d^{-}(v)=\mult_D(v)$: $D_k$ is balanced, so the unitig lemma, the bundle-chain
lemma and the splitting theorem apply to it verbatim. Three consequences. \emph{One:}
multiplicities add across strands, so a family with $a$ forward and $b$ reverse copies has
multiplicity $a+b$ in $D_k$; for \emph{E.\ coli} the longest multiplicity-seven chain has
span $728$\,bp and its seven occurrences are the seven operons, one each, the two
minus-strand ones on the reverse strand of $D_k$. \emph{Two:} $D_k$ is connected exactly
when $S$ contains an inverted repeat of length $\ge k$, so the inverted repeats are
precisely what glues the strands; $\varphi$X174 already qualifies at $k=12$, through
\texttt{ATTAAGCTCATT} at $441$ and its reverse complement at $1448$. \emph{Three:} the
involution $\rho(v)=\bar v$ reverses arcs, hence carries $D_k$ to its opposite; since
$\Delta^{\mathrm{op}}=\Delta^{\mathsf T}$ for a balanced digraph and the two Bowen--Franks
cokernels are dual, $\rho$ \emph{dualises} $\K(D_k)$ instead of splitting it --- which is
what the warning against $\pm1$-eigenspace decompositions was pointing at. The phage's
celebrated single global bit is a single-strand phenomenon: $\K(G_{12})=\Z/2$ but
$\K(D_{12})=\Z/33$. Six checks, all exact. What is \emph{not} done: $D_k$ is the double
cover, not the bidirected quotient, and $\K(D_k)$ does not count double-stranded
reconstructions; that problem is left exactly where it was.
\end{chapterabstract}

\section{Introduction}\label{Bio6:sec:intro}

For a circular sequence $S$ and $k\ge1$, $G_k(S)$ has the distinct $k$-mers as vertices
and the occurrences of $(k{+}1)$-mers as arcs, and
$\K\bigl(G_k(S)\bigr)\cong\bigoplus_c(\Z/r_c)^{\ell_c-1}\oplus\K(\Sk)$ by the bundle-chain
splitting theorem [Ch.~\ref{chap:Bio3}]. The theory is single-strand throughout [Chs.~\ref{chap:Bio1}, \ref{chap:Bio2}, \ref{chap:Bio3}, \ref{chap:Bio4}],
and [Ch.~\ref{chap:Bio5}] measured what that costs on a real genome: of the seven \emph{rrn} operons
of \emph{E.\ coli} K-12, five lie on the $+$ strand and two on the $-$, and since a repeat
of $G_k(S)$ requires co-oriented copies, the operons split into a multiplicity-five family
and a multiplicity-two family, with multiplicity seven occurring nowhere.

The obvious repair is to work with the bidirected graph that assemblers actually build,
whose vertices are the classes $\{v,\bar v\}$; [Chs.~\ref{chap:Bio1}, \ref{chap:Bio2}] left this open and warned
against the tempting shortcut of decomposing $\K$ into $\pm1$-eigenspaces of the
reverse-complement involution, since an involution splits an abelian group only after $2$
is inverted and the $2$-torsion is where the interesting phenomena live. That warning
stands. What this note observes is that most of what one wants from the bidirected picture
is already available, without any new lemma, from a balanced digraph.

\section{The double cover}\label{Bio6:sec:cover}

\begin{definition}[Reverse-complement double cover]\label{Bio6:def:cover}
Let $\bar S$ be the reverse complement of $S$, again circular of length $N$. Let
$D_k=D_k(S)$ be the multidigraph whose vertices are the $k$-mers occurring in $S$ or in
$\bar S$ and whose arcs are the occurrences of $(k{+}1)$-mers in $S$ together with those
in $\bar S$; it has $2N$ arcs. Write $\rho(v)=\bar v$.
\end{definition}

\begin{lemma}[Balance]\label{Bio6:lem:balance}
$\mult_{D}(v)=\mult_S(v)+\mult_S(\bar v)$ and $d^{+}(v)=d^{-}(v)=\mult_D(v)$ for every
vertex. Hence $D_k$ is a balanced multidigraph.
\end{lemma}

\begin{proof}
$\mult_{\bar S}(v)=\mult_S(\bar v)$, because a copy of $v$ in $\bar S$ is a copy of
$\bar v$ in $S$; the first formula follows. In each of the two circular sequences
separately, every occurrence of $v$ contributes one arc out of $v$ and one arc into $v$,
as in [Ch.~\ref{chap:Bio1}, Lem.~4]; summing the two gives the second.
\end{proof}

\begin{corollary}[Nothing new is needed]\label{Bio6:cor:applies}
The unitig contraction lemma [Ch.~\ref{chap:Bio1}, Lem.~4], the bundle-chain lemma
[Ch.~\ref{chap:Bio3}, Lem.~2], the canonicity of the skeleton [Ch.~\ref{chap:Bio3}, Lem.~6] and the splitting
theorem [Ch.~\ref{chap:Bio3}, Thm.~5] apply to $D_k$ verbatim, and a vertex of $D_k$ survives unitig
contraction if and only if $\mult_S(v)+\mult_S(\bar v)\ge2$ [Ch.~\ref{chap:Bio5}, Lem.~2.1].
\end{corollary}

That last clause is the practical point: the branch vertices of $D_k$ are the $k$-mers
that are repeated \emph{when both strands are counted}, so the entire fast pipeline of
[Ch.~\ref{chap:Bio5}] runs on $D_k$ unchanged, with the positions of $\bar S$ appended to those of
$S$.

\begin{proposition}[The strands are glued by the inverted repeats]\label{Bio6:prop:conn}
$D_k$ is connected if and only if some $k$-mer $w$ occurs in $S$ together with $\bar w$;
otherwise $D_k=G_k(S)\sqcup G_k(\bar S)$, two isomorphic components.
\end{proposition}

\begin{proof}
The arcs coming from $S$ form a closed walk through every $k$-mer of $S$, so those
vertices lie in one component, and likewise for $\bar S$; the two vertex sets are the
$k$-mers of $S$ and their reverse complements. They meet exactly when some $w$ and
$\bar w$ both occur in $S$.
\end{proof}

So the double cover is interesting precisely to the extent that the genome has inverted
repeats of length at least $k$. Both of our genomes qualify: $\varphi$X174 at $k=12$
through \texttt{ATTAAGCTCATT} at position $441$, whose reverse complement sits at $1448$;
\emph{E.\ coli} at $k=150$ through a $150$-mer at $19\,796$ whose reverse complement is at
$3\,584\,046$ (check~(C)).

\section{Multiplicities add across strands}\label{Bio6:sec:merge}

\begin{theorem}[Strand merging]\label{Bio6:thm:merge}
Let $W$ be a sequence occurring $a$ times on the forward strand of $S$ and $b$ times on the
reverse, with $a+b\ge1$ and $W$ not equal to $\bar W$. In $G_k(S)$ the $k$-mers of $W$
have multiplicity $a$ and those of $\bar W$ multiplicity $b$. In $D_k$ both have
multiplicity $a+b$, and the two families are exchanged by $\rho$. If, moreover, the $a+b$
copies agree over a window of $s+k-1$ letters, that window contributes a bundle chain of
multiplicity $a+b$ and length $s$ to $D_k$, hence a summand $(\Z/(a+b))^{\,s-1}$ to
$\K(D_k)$, and another from its $\rho$-image.
\end{theorem}

\begin{proof}
A reverse-strand copy of $W$ is a forward occurrence of $\bar W$, so
$\mult_S(w)=a$ and $\mult_S(\bar w)=b$ for a $k$-mer $w$ of $W$ in general position; by
Lemma~\ref{Bio6:lem:balance}, $\mult_D(w)=a+b$, and symmetrically for $\bar w$. The last
sentence is Lemma~\ref{Bio6:lem:balance} followed by [Ch.~\ref{chap:Bio3}, Lem.~2] applied to the
balanced digraph $D_k$.
\end{proof}

\begin{corollary}[The seven operons]\label{Bio6:cor:rrna}
For \emph{E.\ coli} K-12 the seven \emph{rrn} operons carry, in $D_k$, bundle chains of
multiplicity seven. At $k=150,100,51$ the longest such chain has span $728$\,bp and its
seven occurrences are the seven operons, one each:
\[
\begin{array}{lll}
+\,224\,044\to223\,771, & -\,2\,730\,735\to2\,726\,069, & -\,3\,428\,340\to3\,423\,423,\\
+\,3\,942\,081\to3\,941\,808, & +\,4\,035\,804\to4\,035\,531, & +\,4\,166\,932\to4\,166\,659,\\
+\,4\,208\,420\to4\,208\,147, & &
\end{array}
\]
the sign giving the strand of $D_k$ on which the occurrence lies and the arrow pointing to
the operon it meets. The multiplicity-seven chains carry more $k$-mers than the chains of
any other multiplicity over the operons: $4154$ at $k=150$, $5254$ at $k=100$, $6660$ at
$k=51$ (check~(M)).
\end{corollary}

\begin{remark}[What was wrong with $5+2$, exactly]\label{Bio6:rem:52}
Nothing was wrong with it; it was the correct answer to a single-strand question. The
content of Corollary~\ref{Bio6:cor:rrna} is that the question has a balanced-digraph answer
too, and that the object which answers it is not the bidirected quotient but the double
cover, on which the existing theorems already hold. The residual multiplicities
$2,3,4,5,6$ that still appear over the operons in $D_k$ are the sub-families on which only
some of the seven copies agree --- the operons are not identical to one another, on either
strand.
\end{remark}

\section{The involution dualises, it does not split}\label{Bio6:sec:rho}

\begin{proposition}[$\rho$ and the opposite digraph]\label{Bio6:prop:rho}
$\rho$ is an involution of the vertex set of $D_k$ carrying the arc for a $(k{+}1)$-mer
$x$ to the arc for $\bar x$ with its direction reversed. Hence $\rho$ is an isomorphism
$D_k\to D_k^{\mathrm{op}}$, and it induces an isomorphism
\[
\K(D_k)\;\xrightarrow{\ \sim\ }\;\K(D_k)^{\vee}
\]
onto the Pontryagin dual, i.e.\ a perfect pairing on $\K(D_k)$.
\end{proposition}

\begin{proof}
The prefix of $\bar x$ is the reverse complement of the suffix of $x$ and vice versa, so
the arc of $\bar x$ runs $\rho(\text{suffix})\to\rho(\text{prefix})$: direction reverses,
and $\rho$ is an isomorphism onto the opposite digraph. For a balanced digraph
$\Delta^{\mathrm{op}}=D^{+}-A^{\mathsf T}=\Delta^{\mathsf T}$, so
$\K(D_k^{\mathrm{op}})=\operatorname{Tor}\coker\Delta_{\hat s}$ while
$\K(D_k)=\operatorname{Tor}\coker\Delta_{\hat s}^{\mathsf T}$, and these two are
canonically dual by the Bowen--Franks pairing \cite[Lem.~1.5]{ANCFT8}. Composing gives the
stated isomorphism.
\end{proof}

\begin{remark}[Against the eigenspace decomposition]\label{Bio6:rem:nosplit}
[Ch.~\ref{chap:Bio1}, Open (c)] warned that $\K$ should not be split into $\pm1$-eigenspaces of
$\rho$, because that discards the $2$-torsion. Proposition~\ref{Bio6:prop:rho} is what happens
instead: an involution that reverses arcs cannot act on $\K(D_k)$ at all --- it does not
preserve the presentation --- but it identifies $\K(D_k)$ with its dual. No inversion of
$2$ occurs anywhere, and the $2$-torsion is untouched. We do not determine whether the
resulting pairing is symmetric or alternating; that question, and whether it is
independent of the choice of sink, are left open.
\end{remark}

\section{Computations}\label{Bio6:sec:comp}

\begin{table}[ht]
\centering\footnotesize
\begin{tabular}{rlrrrll}
\toprule
$k$ & graph & branch & chains & skeleton & $\K(\Sk)$ & length-driven\\
\midrule
12 & $G_k$ &   2 &  1 &  1 & $0$ & $\Z/2$\\
   & $D_k$ &   7 &  0 &  7 & $\Z/3\oplus\Z/11$ & $0$\\[2pt]
11 & $G_k$ &  10 &  2 &  8 & $\Z/3\oplus\Z/11$ & $(\Z/2)^{2}$\\
   & $D_k$ &  30 &  7 & 23 & $(\Z/2)^{2}\oplus\Z/49451$ & $(\Z/2)^{7}$\\[2pt]
10 & $G_k$ &  47 &  8 & 37 & $\Z/4\oplus\Z/1076338579$ & $(\Z/2)^{10}$\\
   & $D_k$ & 125 & 23 & 97 & $(\Z/4)^{2}\oplus\Z/79624452132275189516228869$ & $(\Z/2)^{28}$\\
\bottomrule
\end{tabular}
\caption{$\varphi$X174: the single-strand graph and its double cover. All entries exact
(check~(K)); the single-strand rows reproduce [Ch.~\ref{chap:Bio3}, Tab.~1].}
\label{Bio6:tab:phix}
\end{table}

\begin{remark}[The phage's single bit is single-strand]\label{Bio6:rem:bit}
$\K(G_{12})=\Z/2$: the $12$-mer spectrum of $\varphi$X174 leaves exactly two global
reconstructions, and [Ch.~\ref{chap:Bio2}] located the bit as a pair of interleaved $12$-mers. On
the double-stranded spectrum the answer is different: $\K(D_{12})=\Z/3\oplus\Z/11=\Z/33$,
and it is entirely arrangement --- $D_{12}$ has no bundle chain at all, so the
length-driven summand is trivial. The two statements do not conflict; they answer
different questions, and the contrast is a fair measure of how much of the single-strand
picture is an artefact of ignoring one strand.
\end{remark}

\begin{table}[ht]
\centering\small
\begin{tabular}{rrrrrrr}
\toprule
$k$ & branch & chains & skeleton & $\log_{10}$ length-driven & $\log_{10}|\K(\Sk)|$
& $\log_{10}|\K(D_k)|$\\
\midrule
150 & 39\,768 & 278 & 312 & 20\,987.91 & 156.51 & 21\,144.42\\
100 & 44\,808 & 396 & 434 & 23\,394.31 & 208.92 & 23\,603.24\\
 51 & 53\,380 & 798 & 886 & 27\,405.96 & 395.80 & 27\,801.75\\
\bottomrule
\end{tabular}
\caption{\emph{E.\ coli} K-12 MG1655, the double cover. Structural entries exact;
$\log_{10}|\K(\Sk)|$ by float LU. Compare [Ch.~\ref{chap:Bio5}, Tab.~1]: on one strand the same $k$
give skeletons of $115$, $175$ and $366$ vertices and arrangement exponents $49.71$,
$71.11$ and $143.17$.}
\label{Bio6:tab:ecoli}
\end{table}

Two readings of Table~\ref{Bio6:tab:ecoli}. The skeleton roughly triples ($115\to312$ at
$k=150$) and the arrangement exponent does too ($49.71\to156.51$): inverted repeats are a
large part of the irreducible ambiguity of this genome and the single-strand theory sees
none of them. And $\log_{10}|\K(D_k)|$ exceeds twice the single-strand value by $166.6$,
$222.1$ and $368.1$ at $k=150,100,51$ --- a crude but honest measure of what
double-strandedness adds beyond simply having two strands to reconstruct.

\section{Verification}\label{Bio6:sec:battery}

\begin{table}[ht]
\centering\small
\begin{tabular}{clc}
\toprule
key & check & mode\\
\midrule
(B) & Lemma~\ref{Bio6:lem:balance}: $D_k$ balanced, multiplicities add across strands & exact\\
(C) & Proposition~\ref{Bio6:prop:conn}: connected iff an inverted repeat of length $\ge k$, with witness & exact\\
(R) & Proposition~\ref{Bio6:prop:rho}: $\rho$ is an involution with $A(\rho u,\rho v)=A(v,u)$ & exact\\
(M) & Corollary~\ref{Bio6:cor:rrna}: the multiplicity-seven chains and their seven operons & exact\\
(K) & Table~\ref{Bio6:tab:phix}: $\K(D_k)$ exactly, beside $\K(G_k)$ & exact\\
(E) & Table~\ref{Bio6:tab:ecoli} & exact/numeric\\
\bottomrule
\end{tabular}
\caption{The verification battery (\texttt{rc\_double\_cover.py}), $6/6$, about $33$
seconds.}
\label{Bio6:tab:battery}
\end{table}

\section{Scope}\label{Bio6:sec:scope}

\emph{Proved.} Lemma~\ref{Bio6:lem:balance}, Corollary~\ref{Bio6:cor:applies},
Proposition~\ref{Bio6:prop:conn}, Theorem~\ref{Bio6:thm:merge}, Corollary~\ref{Bio6:cor:rrna} and
Proposition~\ref{Bio6:prop:rho}. All are short; the work is in noticing that the double cover
is balanced and that this is enough.

\emph{What this is not.} $D_k$ is the \emph{double cover}, not the bidirected quotient.
The two are different objects and only the first is a digraph. In particular:
\begin{itemize}
\item $\K(D_k)$ does \emph{not} count double-stranded reconstructions. By BEST,
$|\K(D_k)|\prod_v(d^{+}(v)-1)!$ is the number of arc-rooted Eulerian circuits of $D_k$,
and such a circuit traverses all $2N$ arcs in one closed walk, which is not a genome. A
double-stranded circular molecule corresponds instead to a decomposition of $D_k$ into two
$\rho$-conjugate circuits. Counting those is not solved here.
\item No sandpile theory for the bidirected quotient is offered. The tools
\cite{Zaslavsky} named in [Chs.~\ref{chap:Bio1}, \ref{chap:Bio2}] remain the place to start, and the obstruction
noted there --- that reverse complementation reverses arc direction, so the quotient map
is not a covering of digraphs and \cite{RT} does not apply --- is untouched.
\item The pairing of Proposition~\ref{Bio6:prop:rho} is not identified as symmetric or
alternating, and we have not checked that it is independent of the sink.
\end{itemize}

\emph{Numeric.} $\log_{10}|\K(\Sk)|$ in Table~\ref{Bio6:tab:ecoli} is a float LU determinant of
a $311$- to $885$-dimensional matrix; everything else is exact integer work.

\emph{Open.} (a) The two questions just listed. (b) Whether
Theorem~\ref{Bio6:thm:merge} has a converse: is every multiplicity-$r$ family of $D_k$ a
strand-split family of $S$, and can $a$ and $b$ be recovered from $D_k$ alone?
(c) The $26$-digit prime in Table~\ref{Bio6:tab:phix} we do not interpret, as in [Ch.~\ref{chap:Bio3}].

\section*{Provenance of this chapter}

Unrefereed preprint. Every statement above is proved in full and every number is produced
by \texttt{rc\_double\_cover.py} from the two public genomes. The task that prompted this
note asked for a bidirected bundle-chain lemma yielding summands
$(\Z/2r)^{s-1}$; no such lemma is needed for the double cover, and
Theorem~\ref{Bio6:thm:merge} shows the correct multiplicity is $a+b$, the number of copies on
both strands together, not $2r$. The same task asked for ``the algebraic proof that the
seven operons participate in the skeleton as a single seven-copy family''; that is
Corollary~\ref{Bio6:cor:rrna}, with the caveat of Remark~\ref{Bio6:rem:52} --- they do so in the
double cover, and there are two such families exchanged by $\rho$, not one.
